\documentclass[18pt, a3paper, landscape, margin=2mm, innermargin=2mm, blockverticalspace=2mm, colspace=5mm]{tikzposter}

\usepackage[utf8]{inputenc}
\usepackage{amsmath, amssymb, amsfonts, physics}

\usetheme{Basic}
\usecolorstyle{Default}

\colorlet{backgroundcolor}{white}
\colorlet{blocktitlebgcolor}{black!10}
\colorlet{blocktitlefgcolor}{black}
\colorlet{blockbodybgcolor}{white}
\colorlet{blockbodyfgcolor}{black}

\title{\vspace{-1.5cm}\parbox{\linewidth}{\centering \Big 
A Rigorous and Stable Implementation of the \\
Crank--Nicolson Method for the TDSE
}}
\begin{document}
\makeatletter
\setlength{\TP@titleinnersep}{0pt}   % removes inner padding
\setlength{\TP@titletotopverticalspace}{0pt} % removes top gap
\setlength{\TP@titletoblockverticalspace}{1mm} % adjust spacing after title
\makeatother
\maketitle

\begin{columns}
% ==================== COLUMN 1 ====================
\column{0.33}

\block{1. Physical Problem: Time Evolution in Quantum Mechanics}{
The time-dependent Schr\"odinger equation (TDSE) governs the evolution of quantum states:
\begin{equation}
i\hbar \frac{\partial \Psi(x,t)}{\partial t}
=
\left[
-\frac{\hbar^2}{2m}\frac{\partial^2}{\partial x^2} + V(x)
\right]\Psi(x,t)
\end{equation}

\textbf{Key Requirement:} The total probability must be conserved:
\[
\int |\Psi(x,t)|^2 dx = 1
\]

\textbf{Interpretation:}
\begin{itemize}
\item The Hamiltonian $\hat{H}$ controls energy + dynamics
\item Evolution must be \textbf{unitary}
\item Numerical methods must preserve norm (very important!)
\end{itemize}
}

\block{2. Discretization of Space and Time}{
We discretize:
\[
x_j = j\Delta x, \quad t_n = n\Delta t
\]

Second derivative (central difference):
\begin{equation}
\frac{\partial^2 \Psi}{\partial x^2}
\approx
\frac{\Psi_{j+1}^n - 2\Psi_j^n + \Psi_{j-1}^n}{(\Delta x)^2}
\end{equation}

\textbf{Result:} PDE $\rightarrow$ System of algebraic equations

\textbf{Insight:}
\begin{itemize}
\item Smaller $\Delta x$ → higher spatial accuracy
\item Smaller $\Delta t$ → better temporal resolution
\item Tradeoff: accuracy vs computational cost
\end{itemize}
}

\block{3. Why Standard Methods Fail}{
\textbf{FTCS (Explicit Method):}
\begin{itemize}
\item Easy to implement
\item But \textbf{unconditionally unstable}
\item Amplifies numerical noise → blows up
\end{itemize}

\textbf{BTCS (Implicit Method):}
\begin{itemize}
\item Stable
\item But introduces \textbf{artificial damping}
\item Violates quantum probability conservation
\end{itemize}

\textbf{Conclusion:}
We need a method that is:
\begin{itemize}
\item Stable
\item Accurate
\item Norm-preserving
\end{itemize}
}

% ==================== COLUMN 2 ====================
\column{0.34}

\block{4. Crank–Nicolson Scheme}{
The Crank–Nicolson method averages FTCS and BTCS:

\begin{equation}
i\hbar \frac{\Psi^{n+1}_j - \Psi^n_j}{\Delta t}
=
\frac{1}{2}\hat{H}_{FD}\Psi^{n+1}_j
+
\frac{1}{2}\hat{H}_{FD}\Psi^n_j
\end{equation}

Rewriting:
\begin{equation}
\left( I + \frac{i\Delta t}{2\hbar}H \right)\Psi^{n+1}
=
\left( I - \frac{i\Delta t}{2\hbar}H \right)\Psi^n
\end{equation}

\textbf{Key Idea:}
\begin{itemize}
\item Implicit + symmetric in time
\item Second-order accurate in both space and time
\end{itemize}
}

\block{5. Unitarity and Stability}{
Define propagator:
\[
U_{CN} = (I + i\gamma H)^{-1}(I - i\gamma H)
\]

This is a \textbf{Cayley transform} of a Hermitian operator.

\textbf{Result:}
\[
U^\dagger U = I
\]

\textbf{Implications:}
\begin{itemize}
\item Exact probability conservation
\item No artificial damping
\item No instability
\end{itemize}

\textbf{This is why CN is the gold standard for TDSE.}
}

\block{6. Matrix Formulation}{
Let:
\[
\alpha = \frac{i\Delta t}{2\hbar}, \quad
\beta = \frac{\hbar^2}{2m(\Delta x)^2}
\]

Hamiltonian becomes:
\[
H\Psi_j = -\beta\Psi_{j-1} + (2\beta + V_j)\Psi_j - \beta\Psi_{j+1}
\]

We obtain:
\[
A\Psi^{n+1} = B\Psi^n
\]

\textbf{Important:}
\begin{itemize}
\item $A$, $B$ are tridiagonal
\item Only nearest-neighbour coupling
\end{itemize}
}

% ==================== COLUMN 3 ====================
\column{0.34}

\block{7. Efficient Solution: Thomas Algorithm}{
Instead of inverting matrices:

\[
\mathcal{O}(N^3) \rightarrow \mathcal{O}(N)
\]

\textbf{Thomas Algorithm Steps:}
\begin{enumerate}
\item Forward elimination
\item Modify coefficients
\item Back substitution
\end{enumerate}

\textbf{Why it works:}
\begin{itemize}
\item Exploits tridiagonal structure
\item No unnecessary computations
\end{itemize}
}

\block{8. Boundary Conditions}{
\textbf{Dirichlet (Infinite Well):}
\[
\Psi_0 = \Psi_N = 0
\]

\textbf{Periodic:}
\[
\Psi_0 = \Psi_{N-1}
\]

\textbf{Physical Meaning:}
\begin{itemize}
\item Dirichlet → particle confined
\item Periodic → lattice / solid-state physics
\end{itemize}
}

\block{9. Algorithm Summary (Implementation)}{
\textbf{Step-by-step:}

\begin{enumerate}
\item Initialize $\Psi(x,0)$
\item Construct matrices $A$, $B$
\item For each time step:
\begin{itemize}
\item Compute RHS: $B\Psi^n$
\item Solve $A\Psi^{n+1} = RHS$
\item Normalize (optional check)
\end{itemize}
\item Repeat
\end{enumerate}

\textbf{Output:} Time evolution of wavefunction
}

\block{10. Applications}{
\begin{itemize}
\item Quantum tunneling
\item Wave packet propagation
\item Quantum wells and barriers
\end{itemize}

\textbf{Takeaway:}
Crank–Nicolson provides a perfect balance of:
\begin{center}
\textbf{Accuracy + Stability + Physical Correctness}
\end{center}
}

\end{columns}
\end{document}